% LSPV-3 REPORT::EXTEVAL1::2026-10-05T00:53:00-07:00
\documentclass[11pt]{article}
\usepackage[margin=.88in]{geometry}
\usepackage{amsmath,amssymb}
\begin{document}
\section*{LSPV-3 EXTEVAL1: declared SU05/Hadamard seed and evaluation exponents}
\textbf{STATUS: OPEN.} P versus NP remains open.
\section*{Primary-literature result}
vMS Definition 25 defines a local extractor as computable in time polynomial in seed/output lengths. Lemma 26 imports the Reed--Muller extractor with seed length $O(r\log p)$, output length $\Delta$, and only $\operatorname{poly}(\Delta,r,\log p)$ oracle accesses. This is enough for polynomial source-side evaluation but does not expose a numerical leading seed constant.
A later space-efficient analysis of the Shaltiel--Umans extractor explicitly notes that the q-ary extractor can be run while storing the current orbit vector and that the final Hadamard conversion is computable in polynomial time and seed-linear space. This supports streaming evaluation but does not repair SAT acquisition.
\section*{Declared SU05/Hadamard seed}
For one RMV line the active SU orbit lives in $B=\mathbb F_p^d$ with $d=r/2$. The declared q-ary orbit seed is $v\in B$. Hadamard conversion adds one mask/coordinate of $k=\lceil\log_2p\rceil$ bits. Therefore the declared realization has
\[
s_{\rm dec}=d\log_2p+k
=\left(\frac r2+1\right)\log_2p+O(1),
\]
and relative to parent explicit-table size $N=p^r$,
\[
\boxed{\lambda_{\rm dec}=\frac12+\frac1r+o(1).}
\]
This is exact for the declared Hadamard realization up to ceiling/lower-order encoding terms. It is not asserted as the hidden constant of vMS Lemma 26.
\section*{Direct branch}
SEEDBRANCH1 gives exactly $2p^{r/2}=N^{1/2+o(1)}$ distinct MV lines and separates direct extractor work from recursive RMV calls. Thus the direct branch exponent is
\[
D_{\rm ext}=\frac12+\lambda_{\rm dec}+g_{\rm ext}
=1+\frac1r+g_{\rm ext}+o(1).
\]
\section*{Evaluation cost: random-access table}
If the line oracle is already supplied as a random-access evaluation table, one orbit query plus Hadamard projection uses only polylogarithmic coordinate arithmetic. For a fixed output coordinate this gives
\[
\boxed{g_{\rm ext}=0}
\]
in exponent-of-$N$ accounting. Generating all output coordinates contributes the declared output-length factor separately if they are all materialized.
This is a representation result, not an acquisition result.
\section*{Evaluation cost: dense coefficient representation}
For a line polynomial on $d=r/2$ variables of total degree at most $\Delta$, the number of dense coefficients is
\[
M_d=\binom{\Delta+d}{d}\le (\Delta+d)^d.
\]
The vMS parameter regime has $p\ge\Delta^{100}$, hence $\Delta\le p^{1/100}$. Ignoring lower-order $d$ factors when $d\ll\Delta$,
\[
M_d\le p^{d/100+o(d)}
=N^{1/200+o(1)}.
\]
A straightforward dense evaluation therefore has the conservative representation-side candidate
\[
\boxed{g_{\rm coeff}\le\frac1{200}+o(1)}
\]
per field evaluation, provided the coefficient representation has already been acquired. This bound is deliberately conservative; Horner/tensor structure may improve constants, but acquisition of coefficients from the canonical proof table is separately charged.
The RMV output prefix has $m\le\Delta^{1/100}$, so materializing all $m$ orbit symbols adds at most
\[
\frac{1}{10000r}
\]
to the parent-$N$ exponent. This does not threaten the subquadratic threshold.
\section*{Combined declared recurrence}
Using the dense-coefficient candidate gives
\[
C_i\ge
\max\left\{
r_i,\;
1+\frac1r+\frac1{200}+\frac1{10000r}+o(1),\;
\frac12+\frac12C_j
\right\}.
\]
At the smallest recursive source dimension $r=4$, the direct exponent is at most approximately
\[
1+\frac14+\frac1{200}+\frac1{40000}<1.256,
\]
under the already-acquired dense-coefficient representation. At the terminal $r=2$ leaf there is no recursive child; the analogous direct exponent is approximately below $1.506$.
Thus the declared source-side generator is comfortably subquadratic in this explicit-table normalization after representation acquisition.
\section*{Critical interpretation}
This does \emph{not} produce a SAT speedup or lower bound. The source object itself is already an encoded/LDE representation. Constructing that representation from an arbitrary SAT computation, proving that an arbitrary SAT decider must expose it, or proving a SAME\_REALIZATION return into SAT remain exactly the missing acquisition/transfer theorems. Moving LDE construction into preprocessing would merely hide the charge unless its source is separately bounded.
\section*{Literature comparison}
The general SU05 theorem advertises seed length $(1+\alpha)\log n$ for suitable extractor parameterizations, but that global source-length normalization is not the same as the present parent-table normalization and must not be substituted for $\lambda_{\rm dec}$. Later algebraic extractor work such as Guruswami--Umans--Vadhan changes the extractor architecture; it may improve seed parameters but does not automatically preserve the RMV reconstruction/SAME\_REALIZATION interface. It is therefore literature context, not a drop-in replacement.
\section*{Falsification and surviving target}
FALSIFIED: the concern that local extractor evaluation necessarily destroys the square-root recursion once the exact declared SU05/Hadamard evaluator is used on an already available line representation. The source-side evaluator can be polynomially local/streaming, and explicit representations give small exponent charges.
NOT FALSIFIED: acquisition. The arbitrary-algorithm bridge is now even more clearly the dominant missing theorem.
\section*{Next target}
Charge acquisition of one canonical LDE line representation from the perfectly accepted PCPP proof, distinguishing random-access table, dense coefficients, and implicit tensor-grid evaluation. Determine whether one shared parent representation can service all $2p^{r/2}$ line children without paying a fresh $N$-scale construction per line. Then apply the SAME\_REALIZATION hostile wrappers. If acquisition is $N^{1+o(1)}$ once globally and line access is sublinear thereafter, the source-side recurrence remains favorable; if acquisition repeats per child, the gain disappears.
\section*{Status}
PROVED/EXACT for the declared realization: seed decomposition and $\lambda_{\rm dec}=1/2+1/r+o(1)$. PROVED source fact: vMS local extractor has polynomial evaluation/query complexity. CONDITIONAL representation accounting: $g=0$ for preexisting random-access table and $g\le1/200+o(1)$ for an already-acquired dense coefficient representation under the stated parameter inequalities. OPEN: acquisition, SAT return, nonrelativization, nonalgebrization, and P versus NP.
\begin{thebibliography}{4}
\bibitem{SU05} R. Shaltiel and C. Umans, \emph{Simple Extractors for All Min-Entropies and a New Pseudorandom Generator}, J. ACM 52(2):172--216, 2005.
\bibitem{vMS25} D. van Melkebeek and N. Mocelin Sdroievski, \emph{Instance-Wise Hardness and Refutation versus Derandomization for Arthur--Merlin Protocols}, Computational Complexity 34:15, 2025.
\bibitem{GUV09} V. Guruswami, C. Umans, and S. Vadhan, \emph{Unbalanced Expanders and Randomness Extractors from Parvaresh--Vardy Codes}, J. ACM 56(4), 2009.
\bibitem{TV19} S. Tell and R. van Melkebeek, \emph{Typically-Correct Derandomization for Small Time and Space}, CCC 2019.
\end{thebibliography}
% END_OF_VERIFIED_REPORT::LSPV-3::EXTEVAL1
\end{document}
